% =========================================================================
\section{Seminarios prácticos}\label{sec:seminarios}
% =========================================================================
Esta sección reúne los ejercicios del tema resueltos con todo detalle. Cada
ejercicio sigue la misma estructura: un \emph{enunciado}, la \emph{idea} que
guía la resolución (qué estamos buscando y por qué el método funciona), la
\emph{resolución} paso a paso y, cuando merece la pena, un breve comentario con
la \emph{interpretación} geométrica o física del resultado. Los seminarios
\ref{sem:pendiente}--\ref{sem:levi} trabajan con operadores diferenciales;
los demás verifican, con campos concretos, los tres teoremas integrales de la
\cref{sec:integral}.

% =========================================================================
\subsection{Seminario 1. Campos escalares: la pendiente máxima de una campana}\label{sem:pendiente}
% =========================================================================
\begin{ejercicio}[Pendiente máxima]\label{ejer:pendiente}
Para el campo escalar $f(\vc{r})=\eul^{-(x^{2}+y^{2})}$, calcular a qué distancia
del origen la pendiente es máxima.
\end{ejercicio}

\begin{resolucion}
\textbf{Idea.} Imaginamos $f:\R^{2}\to\R$ como una superficie $z=f(x,y)$: una
campana centrada en el origen. En un punto $P=(x,y)$ podemos movernos en
cualquier dirección del plano y la derivada direccional nos dice cuál es la
pendiente de la superficie en esa dirección. Por la \cref{prop:maximo}, la
mayor de todas esas pendientes en $P$ es exactamente $\norm{\nab f(P)}$, y se
alcanza al ir en la dirección del gradiente. La pregunta se transforma
entonces en: \emph{¿dónde es máximo $\norm{\nab f}$?}

\textbf{Paso 1: el gradiente.} Por la regla de la cadena
$[g(f(x))]'=g'(f(x))\,f'(x)$,
\[
  \pd{f}{x}=-2x\,\eul^{-(x^{2}+y^{2})},\qquad
  \pd{f}{y}=-2y\,\eul^{-(x^{2}+y^{2})},
  \quad\Longrightarrow\quad
  \nab f=-2\,\eul^{-(x^{2}+y^{2})}\,(x\,\ux+y\,\uy),
\]
cuyo módulo es $\norm{\nab f}=2\,\eul^{-(x^{2}+y^{2})}\sqrt{x^{2}+y^{2}}$.

\textbf{Paso 2: aprovechar la simetría.} El módulo depende solo de la distancia
al origen. Pasando a polares, $x=r\cos\varphi$, $y=r\sen\varphi$, con
$r=\sqrt{x^{2}+y^{2}}$ ($x^{2}+y^{2}=r^{2}$), queda
$\norm{\nab f(r,\varphi)}=\norm{\nab f(r)}=2r\,\eul^{-r^{2}}$. El problema
bidimensional se ha convertido en uno unidimensional: maximizar
\[
  g(r)=2r\,\eul^{-r^{2}},\qquad r\ge0 .
\]

\textbf{Paso 3: puntos críticos.} Derivamos:
$g'(r)=2\eul^{-r^{2}}+2r(-2r)\eul^{-r^{2}}=2\,\eul^{-r^{2}}(1-2r^{2})$. En un punto
crítico $g'(r)=0$; como $\eul^{-r^{2}}>0$ para todo $r$, necesariamente
$1-2r^{2}=0$, es decir,
\[
  \boxed{\,r=\frac{1}{\sqrt{2}}\approx0{,}707\,}
\]
(en unidades arbitrarias del problema).

\textbf{Paso 4: comprobar que es un máximo.} La segunda derivada es
$g''(r)=4r\,\eul^{-r^{2}}(2r^{2}-3)$ (se obtiene derivando
$2\eul^{-r^{2}}(1-2r^{2})$ por el producto). En $r=1/\sqrt{2}$,
$g''(1/\sqrt{2})=4\tfrac{1}{\sqrt{2}}\,\eul^{-1/2}(1-3)<0$, así que efectivamente
es un máximo, de valor $g(1/\sqrt{2})=\sqrt{2}\,\eul^{-1/2}\approx0{,}858$.
\end{resolucion}

\begin{observacion}[Interpretación]
El resultado es razonable. En el origen $\nab f(\vc{0})=\vc{0}$: la cima de la
campana es plana, la pendiente es nula. Muy lejos, $\eul^{-r^{2}}\to0$ y la
superficie vuelve a aplanarse. Por tanto, tiene que haber una zona intermedia
donde la pendiente sea máxima: crece desde $0$, alcanza su máximo en
$r=1/\sqrt{2}$ y vuelve a decrecer hasta anularse (\cref{fig:sem-pendiente}).
Es, además, el \emph{punto de inflexión} del perfil $\eul^{-r^{2}}$: allí
$f''(r)=(4r^{2}-2)\eul^{-r^{2}}=0$, como corresponde a un punto donde la
pendiente $f'$ es extremal.
\end{observacion}

% Figura: pendiente máxima de la campana gaussiana
\begin{figure}[ht]
\centering
\begin{minipage}[t]{0.52\textwidth}\centering
\begin{tikzpicture}[baseline=(current bounding box.north)]
  \begin{axis}[tema, width=\textwidth, height=5.8cm, xlabel={$r$}, xmin=0, xmax=3, ymin=0, ymax=1.12,
      samples=120, domain=0:3, xtick={0,0.5,0.707,1,1.5,2,3}, xticklabels={0,0{,}5,$\tfrac{1}{\sqrt2}$,1,1{,}5,2,3},
      legend style={at={(0.98,0.97)}, anchor=north east}]
    \addplot[azulacero, line width=1.1pt] {exp(-x^2)};
    \addplot[rojoteo, line width=1.1pt] {2*x*exp(-x^2)};
    \addplot[grispizarra, densely dashed, line width=0.6pt] coordinates {(0.7071,0) (0.7071,0.858)};
    \addplot[only marks, mark=*, mark size=2pt, rojoteo] coordinates {(0.7071,0.8578)};
    \node[font=\scriptsize, rojoteo, anchor=west] at (axis cs:0.78,0.98) {máximo: $\sqrt2\,\eul^{-1/2}\approx0{,}858$};
    \legend{$f(r)=\eul^{-r^{2}}$, $\norm{\nab f}=g(r)=2r\,\eul^{-r^{2}}$}
  \end{axis}
\end{tikzpicture}\\[2pt]{\small\scshape (a)}
\end{minipage}\hfill
\begin{minipage}[t]{0.44\textwidth}\centering
\begin{tikzpicture}[x=0.78cm, y=0.78cm, baseline=(current bounding box.north)]
  \mapaGauss{%
    \draw[white, dashed, line width=1.1pt] (0,0) circle (0.7071);
    \foreach \a in {0,45,...,315}{%
      \draw[-{Stealth[length=4pt]}, line width=0.9pt, white] ({0.7071*cos(\a)},{0.7071*sin(\a)}) -- ++({-0.62*cos(\a)},{-0.62*sin(\a)});}
    \node[white, font=\scriptsize, fill=cm0, fill opacity=0.6, text opacity=1, inner sep=1pt] at (1.45,1.45) {$r=\tfrac{1}{\sqrt2}$};}
\end{tikzpicture}\\[2pt]{\small\scshape (b)}
\end{minipage}
\caption{Pendiente máxima de $f=\eul^{-(x^{2}+y^{2})}$. (a)~Perfil radial $f(r)$ y
módulo de su gradiente $g(r)=\norm{\nab f}$, que crece desde cero en el origen,
alcanza su máximo en $r=1/\sqrt2$ y decae. (b)~Sobre el mapa de calor, la
circunferencia $r=1/\sqrt2$ (discontinua) es el lugar geométrico de máxima
pendiente; las flechas son $\nab f$ allí.}
\label{fig:sem-pendiente}
\end{figure}


% =========================================================================
\subsection{Seminario 2. Gradiente y divergencia: cartesianas frente a esféricas}\label{sem:inverso}
% =========================================================================
\begin{ejercicio}[El campo del inverso del cuadrado]\label{ejer:inverso}
Sea $f(\vc{r})=1/r=(x^{2}+y^{2}+z^{2})^{-1/2}$ (a)~Calcular $\nab f$ en
coordenadas cartesianas y en esféricas. (b)~Para el campo
$\vc{A}(\vc{r})=\ur/r^{2}$ calcular su divergencia en ambos sistemas.
\end{ejercicio}

\begin{resolucion}
\textbf{(a) Gradiente.} En cartesianas, $\partial_x(x^{2}+y^{2}+z^{2})^{-1/2}=-\tfrac12(x^{2}+y^{2}+z^{2})^{-3/2}\cdot2x=-x/r^{3}$, y lo mismo para $y,z$:
\[
  \nab f=\pd{f}{x}\,\ux+\pd{f}{y}\,\uy+\pd{f}{z}\,\uz
  =-\frac{1}{r^{3}}\bigl(x\,\ux+y\,\uy+z\,\uz\bigr)=-\frac{\vc{r}}{r^{3}}=-\frac{\ur}{r^{2}},
\]
donde hemos usado $\vc{r}=r\,\ur$. En esféricas, como $f=f(r)$ solo depende de
$r$, la fórmula del gradiente (\cref{eq:ops-esf}) conserva un único término:
$\nab f=\pd{f}{r}\,\ur=-\dfrac{1}{r^{2}}\,\ur$. Mismo resultado, pero obtenido
sin cálculos: ¡la simetría esférica hace que las esféricas sean el sistema
natural!

\textbf{(b) Divergencia.} El campo $\vc{A}=\ur/r^{2}=\vc{r}/r^{3}$ tiene
componentes cartesianas
$\vc{A}=\dfrac{x\,\ux+y\,\uy+z\,\uz}{(x^{2}+y^{2}+z^{2})^{3/2}}$ (es, salvo signo,
el campo $\nab f$ anterior). En cartesianas,
\[
  \nab\cdot\vc{A}=\pd{A_x}{x}+\pd{A_y}{y}+\pd{A_z}{z},
\]
y cada sumando se calcula con la regla del cociente:
\[
  \pd{A_x}{x}=\pd{}{x}\Bigl[\frac{x}{(x^{2}+y^{2}+z^{2})^{3/2}}\Bigr]
  =\frac{(x^{2}+y^{2}+z^{2})^{3/2}-x\cdot\frac{3}{2}(x^{2}+y^{2}+z^{2})^{1/2}\cdot2x}{(x^{2}+y^{2}+z^{2})^{3}}
  =\frac{1}{r^{3}}-\frac{3x^{2}}{r^{5}} .
\]
Sumando las tres,
\[
  \nab\cdot\vc{A}=\frac{3}{r^{3}}-\frac{3\,(x^{2}+y^{2}+z^{2})}{r^{5}}=\frac{3}{r^{3}}-\frac{3}{r^{3}}=0 .
\]
En esféricas, con $A_r=1/r^{2}$ y $A_\theta=A_\varphi=0$,
\[
  \nab\cdot\vc{A}=\frac{1}{r^{2}}\pd{}{r}\Bigl(r^{2}\cdot\frac{1}{r^{2}}\Bigr)=\frac{1}{r^{2}}\pd{(1)}{r}=0 .
\]
\end{resolucion}

\begin{observacion}[La singularidad del origen]
La divergencia es nula \emph{en todo punto excepto en el origen}, donde el
campo no está definido (diverge) y la divergencia tiene una contribución
singular: una delta de Dirac. El hecho de que, a pesar de ello, el flujo por
cualquier esfera centrada en el origen valga $4\pi$ (\cref{ej:flujo-r3}) y
\emph{no dependa de su radio} es la huella de esa fuente puntual, como muestra
la \cref{fig:sem-inverso} y se formaliza en el \cref{teo:delta}.
\end{observacion}

% Figura: campo del inverso del cuadrado y flujo constante
\begin{figure}[ht]
\centering
\begin{minipage}[t]{0.47\textwidth}\centering
\begin{tikzpicture}[x=0.84cm, y=0.84cm, baseline=(current bounding box.north)]
  \marcoD{-2.7}{2.7}{-2.7}{2.7}{-2,-1,0,1,2}{-2,-1,0,1,2}
  \ejesD{-2.7}{2.7}{-2.7}{2.7}
  \clipD{-2.7}{2.7}{-2.7}{2.7}{\flechas{\x/((\x*\x+\y*\y)^1.5)}{\y/((\x*\x+\y*\y)^1.5)}%
     {-2.4,-1.8,-1.2,-0.6,0.6,1.2,1.8,2.4}{-2.4,-1.8,-1.2,-0.6,0.6,1.2,1.8,2.4}{0.42}{1.4}}
  \fill[azulprof] (0,0) circle (2pt);
\end{tikzpicture}\\[2pt]{\small\scshape (a)}\ $\vc{A}=\ur/r^{2}$
\end{minipage}\hfill
\begin{minipage}[t]{0.50\textwidth}\centering
\begin{tikzpicture}[baseline=(current bounding box.north)]
  \begin{axis}[tema, width=\textwidth, height=5.4cm, xlabel={$r$}, xmin=0, xmax=3, ymin=0, ymax=22,
      samples=150, domain=0.26:3, legend style={at={(0.97,0.97)}, anchor=north east}]
    \addplot[azulacero, line width=1.1pt] {1/x^2};
    \addplot[rojoteo, line width=1.1pt, domain=0.26:3] {4*pi};
    \legend{$\norm{\vc{A}}=1/r^{2}$, $\Phi(r)=\oiint_{S_r}\vc{A}\cdot\dS=4\pi$}
  \end{axis}
\end{tikzpicture}\\[2pt]{\small\scshape (b)}
\end{minipage}
\caption{El campo $\vc{A}=\ur/r^{2}$. (a)~Flechas radiales cuyo módulo cae como
$1/r^{2}$ al alejarse de la fuente (el origen). (b)~El módulo decrece, pero el
flujo por una esfera de radio $r$, $\Phi=\norm{\vc{A}}\cdot4\pi r^{2}$, es constante
e igual a $4\pi$: el campo «se diluye» exactamente al ritmo en que crece la
superficie, de modo que no hay fuentes salvo la del origen.}
\label{fig:sem-inverso}
\end{figure}


% =========================================================================
\subsection{Seminario 3. Cálculo con índices: el tensor de Levi-Civita}\label{sem:levi}\label{sec:levi-civita}
% =========================================================================
Varias identidades vectoriales de la \cref{prop:productos} y de la
\cref{prop:fundamentales} son tediosas de comprobar componente a componente.
La notación con índices las convierte en cálculos cortos y mecánicos. Para ello
necesitamos una herramienta.

\begin{definicion}[Tensor de Levi-Civita]\label{def:levi-civita}
El \emph{tensor de Levi-Civita} (o símbolo de permutación) en tres dimensiones es
\[
  \varepsilon_{ijk}=\begin{cases}
    +1 & \text{si }(i,j,k)\text{ es una permutación par de }(1,2,3),\\
    -1 & \text{si }(i,j,k)\text{ es una permutación impar de }(1,2,3),\\
    \ \ 0 & \text{si dos índices cualesquiera son iguales.}
  \end{cases}
\]
Es decir, $\varepsilon_{123}=\varepsilon_{231}=\varepsilon_{312}=+1$ (permutaciones
cíclicas) y $\varepsilon_{132}=\varepsilon_{213}=\varepsilon_{321}=-1$.
\end{definicion}

% Figura: tensor de Levi-Civita como orientación cíclica
\begin{figure}[ht]
\centering
\begin{tikzpicture}[>=Stealth, x=1cm, y=1cm]
  \foreach \i/\ang in {1/90,2/-30,3/210}{%
    \node[circle, draw=azulprof, fill=azulclaro!50, minimum size=0.8cm, font=\large, text=azulprof] (n\i) at (\ang:1.5) {\i};}
  \draw[->, line width=1.4pt, azulprof] (n1) to[bend left=18] (n2);
  \draw[->, line width=1.4pt, azulprof] (n2) to[bend left=18] (n3);
  \draw[->, line width=1.4pt, azulprof] (n3) to[bend left=18] (n1);
  \node[font=\small, text=azulprof, align=center] at (0,0) {sentido\\ cíclico};
  \node[anchor=west, align=left, font=\small] at (2.6,0.55)
     {$\textcolor{azulprof}{\varepsilon_{123}=\varepsilon_{231}=\varepsilon_{312}=+1}$\quad(a favor de las flechas)};
  \node[anchor=west, align=left, font=\small] at (2.6,-0.15)
     {$\textcolor{rojoteo}{\varepsilon_{132}=\varepsilon_{213}=\varepsilon_{321}=-1}$\quad(en contra)};
  \node[anchor=west, align=left, font=\small] at (2.6,-0.85)
     {$\textcolor{grispizarra}{\varepsilon_{112}=\varepsilon_{233}=\varepsilon_{121}=\cdots=0}$\quad(índices repetidos)};
\end{tikzpicture}
\caption{El tensor de Levi-Civita codifica la orientación: es $+1$ si los
índices recorren $1,2,3$ en el sentido cíclico de las flechas, $-1$ si lo hacen
en sentido contrario y $0$ si algún índice se repite. Así, $\vc{A}\times\vc{B}$
con $(\vc{A}\times\vc{B})_i=\varepsilon_{ijk}A_jB_k$ reproduce
$\ux\times\uy=\uz$ y los demás productos de la base.}
\label{fig:levi}
\end{figure}


Con el convenio de sumación de Einstein (\cref{not:convenios}), el producto
vectorial, el rotacional, la divergencia y el gradiente se escriben
\begin{equation}\label{eq:indices}
  (\vc{A}\times\vc{B})_i=\varepsilon_{ijk}A_jB_k,\qquad
  (\nab\times\vc{A})_i=\varepsilon_{ijk}\partial_jA_k,\qquad
  \nab\cdot\vc{A}=\partial_iA_i,\qquad
  (\nab f)_i=\partial_if ,
\end{equation}
donde $\partial_i\equiv\partial/\partial x_i$ (derivada \emph{espacial}). En la
primera expresión, $c_i=\sum_{j,k=1}^{3}\varepsilon_{ijk}A_jB_k$ es la componente
$i$-ésima, y los índices repetidos $j,k$ se suman.

\begin{lema}[Propiedades del tensor de Levi-Civita]\label{lema:levi}
\begin{enumerate}[label=(\alph*)]
  \item \emph{Antisimetría:} intercambiar dos índices cualesquiera cambia el
        signo, por ejemplo $\varepsilon_{ijk}=-\varepsilon_{jik}=-\varepsilon_{ikj}$.
  \item \emph{Invariancia cíclica:} $\varepsilon_{ijk}=\varepsilon_{jki}=\varepsilon_{kij}$.
  \item \emph{Constancia:} los valores de $\varepsilon_{ijk}$ son números fijos, así
        que $\partial_l\varepsilon_{ijk}=0$ y puede sacarse de cualquier derivada.
\end{enumerate}
\end{lema}
\begin{proof}
(a) Intercambiar dos índices convierte una permutación par en impar y viceversa
(una transposición cambia la paridad), y deja en cero las de índices repetidos.
(b) Una permutación cíclica de tres elementos es par, luego no cambia el signo.
(c) Es constante por definición.
\end{proof}

\begin{lema}[Identidad $\varepsilon$--$\delta$]\label{lema:eps-delta}
Con $\delta_{ij}$ la delta de Kronecker,
\[
  \varepsilon_{ijk}\,\varepsilon_{ilm}=\delta_{jl}\,\delta_{km}-\delta_{jm}\,\delta_{kl}.
\]
\end{lema}
\begin{proof}
Para $j,k$ fijos, el producto $\varepsilon_{ijk}\varepsilon_{ilm}$ (sumado en $i$) solo
es no nulo si $j\neq k$ y si $\{l,m\}=\{j,k\}$, pues entonces $i$ es el tercer
valor distinto. Si $(l,m)=(j,k)$ ambos factores son iguales y el producto es
$(\pm1)^{2}=+1$; si $(l,m)=(k,j)$ difieren en un intercambio y el producto es
$-1$. Eso es exactamente lo que valen $\delta_{jl}\delta_{km}-\delta_{jm}\delta_{kl}$.
\end{proof}

\begin{lema}[Contracción de un tensor simétrico con otro antisimétrico]\label{lema:sim-anti}
Sean $S_{ij}$ simétrico ($S_{ij}=S_{ji}$) y $T_{ij}$ antisimétrico
($T_{ij}=-T_{ji}$). Entonces $C\equiv S_{ij}\,T_{ij}=\sum_{i,j}S_{ij}T_{ij}=0$.
\end{lema}
\begin{proof}
Los índices sumados son «mudos»: podemos renombrarlos libremente sin alterar la
suma ($i\leftrightarrow j$). Entonces
$C=S_{ji}T_{ji}$. Por simetría y antisimetría, $S_{ji}=S_{ij}$ y $T_{ji}=-T_{ij}$, luego
$C=-S_{ij}T_{ij}=-C$, y por tanto $C=0$.
\end{proof}

Además, para campos de clase $C^{2}$ las derivadas parciales conmutan
(teorema de Schwarz), de modo que el operador $\partial_i\partial_j$ es
\emph{simétrico} en $(i,j)$.

\begin{ejercicio}[Divergencia de un producto vectorial]\label{ejer:div-prod-vect}
Demostrar que $\nab\cdot(\vc{A}\times\vc{B})=\vc{B}\cdot(\nab\times\vc{A})-\vc{A}\cdot(\nab\times\vc{B})$.
\end{ejercicio}
\begin{resolucion}
\textbf{Idea.} Escribimos todo con índices, aplicamos la regla del producto a la
derivada (el símbolo $\varepsilon$ es constante) y reconocemos, en cada término,
la componente de un rotacional.

Por la \cref{eq:indices}, la divergencia es $\partial_i(\cdot)_i$ aplicada a
$(\vc{A}\times\vc{B})_i=\varepsilon_{ijk}A_jB_k$:
\[
  \nab\cdot(\vc{A}\times\vc{B})=\partial_i\bigl(\varepsilon_{ijk}A_jB_k\bigr)
  =\varepsilon_{ijk}\bigl[(\partial_iA_j)\,B_k+A_j\,(\partial_iB_k)\bigr].
\]
\emph{Primer sumando.} Por la invariancia cíclica, $\varepsilon_{ijk}=\varepsilon_{kij}$, así que
$\varepsilon_{ijk}\,\partial_iA_j=\varepsilon_{kij}\,\partial_iA_j=(\nab\times\vc{A})_k$
(la definición del rotacional, con $k$ como índice libre). Por tanto
$\varepsilon_{ijk}(\partial_iA_j)B_k=(\nab\times\vc{A})_kB_k=\vc{B}\cdot(\nab\times\vc{A})$.

\emph{Segundo sumando.} Usamos antisimetría para reordenar:
$\varepsilon_{ijk}=-\varepsilon_{jik}$, luego
$\varepsilon_{ijk}\,\partial_iB_k=-\varepsilon_{jik}\,\partial_iB_k=-(\nab\times\vc{B})_j$, y
$\varepsilon_{ijk}A_j(\partial_iB_k)=-A_j(\nab\times\vc{B})_j=-\vc{A}\cdot(\nab\times\vc{B})$.

Sumando ambos resultados se obtiene la identidad pedida.
\end{resolucion}

\begin{ejercicio}[Rotacional de un escalar por un vector]\label{ejer:rot-fA}
Demostrar que $\nab\times(f\vc{A})=\nab f\times\vc{A}+f\,\nab\times\vc{A}$.
\end{ejercicio}
\begin{resolucion}
Por la definición del rotacional en notación indicial y la regla del producto,
\begin{align*}
  \bigl[\nab\times(f\vc{A})\bigr]_i&=\varepsilon_{ijk}\,\partial_j(fA_k)
  =\varepsilon_{ijk}\bigl[(\partial_jf)\,A_k+f\,(\partial_jA_k)\bigr]\\
  &=\varepsilon_{ijk}(\partial_jf)\,A_k+f\,\varepsilon_{ijk}\,\partial_jA_k
  =(\nab f\times\vc{A})_i+f\,(\nab\times\vc{A})_i,
\end{align*}
donde en el último paso reconocemos el producto vectorial
$(\vc{u}\times\vc{A})_i=\varepsilon_{ijk}u_jA_k$ con $u_j=\partial_jf=(\nab f)_j$ y el
rotacional $(\nab\times\vc{A})_i=\varepsilon_{ijk}\partial_jA_k$. Al valer para cada
componente $i$, vale para el vector.
\end{resolucion}

\begin{ejercicio}[Divergencia de un rotacional]\label{ejer:div-rot}
Demostrar que $\nab\cdot(\nab\times\vc{A})=0$ para todo campo $\vc{A}\in C^{2}$.
\end{ejercicio}
\begin{resolucion}
\textbf{Idea.} Es la contracción de algo simétrico con algo antisimétrico.

En notación indicial, $\nab\cdot(\nab\times\vc{A})=\partial_i\bigl(\varepsilon_{ijk}\partial_jA_k\bigr)$
(la expresión entre paréntesis es la componente $i$-ésima del rotacional).
Como $\varepsilon_{ijk}$ es constante (\cref{lema:levi}), sale de la derivada:
\[
  \nab\cdot(\nab\times\vc{A})=\varepsilon_{ijk}\,\partial_i\partial_jA_k .
\]
Para cada $k$ fijo, $T_{ij}=\varepsilon_{ijk}$ es antisimétrico en $(i,j)$
(\cref{lema:levi}), mientras que $S_{ij}=\partial_i\partial_jA_k$ es simétrico
en $(i,j)$ porque $\vc{A}\in C^{2}$ (teorema de Schwarz): $\partial_i\partial_j=\partial_j\partial_i$.
Por el \cref{lema:sim-anti}, la contracción $S_{ij}T_{ij}$ es nula para cada $k$, y
al sumar en $k$ sigue siendo nula:
\[
  \nab\cdot(\nab\times\vc{A})=\varepsilon_{ijk}\,\partial_i\partial_jA_k=0 .
\]
\end{resolucion}

La misma técnica prueba la otra relación fundamental, y la identidad del doble
rotacional.

\begin{ejercicio}[Rotacional de un gradiente]\label{ejer:rot-grad}
Demostrar que $\nab\times(\nab f)=\vc{0}$ para todo $f\in C^{2}$.
\end{ejercicio}
\begin{resolucion}
$[\nab\times(\nab f)]_i=\varepsilon_{ijk}\,\partial_j(\partial_kf)=\varepsilon_{ijk}\,\partial_j\partial_kf$:
antisimétrico en $(j,k)$ contra simétrico en $(j,k)$ (Schwarz). Por el
\cref{lema:sim-anti}, es cero para cada $i$.
\end{resolucion}

\begin{ejercicio}[Rotacional del rotacional]\label{ejer:rotrot}
Demostrar que $\nab\times(\nab\times\vc{A})=\nab(\nab\cdot\vc{A})-\Lap\vc{A}$.
\end{ejercicio}
\begin{resolucion}
Aplicando dos veces la definición del rotacional,
\[
  [\nab\times(\nab\times\vc{A})]_i=\varepsilon_{ijk}\,\partial_j\bigl(\varepsilon_{klm}\,\partial_lA_m\bigr)
  =\varepsilon_{ijk}\,\varepsilon_{klm}\,\partial_j\partial_lA_m .
\]
Por la invariancia cíclica, $\varepsilon_{ijk}=\varepsilon_{kij}$, y la identidad $\varepsilon$--$\delta$
(\cref{lema:eps-delta}, con el índice contraído $k$ en primera posición),
\[
  \varepsilon_{kij}\,\varepsilon_{klm}=\delta_{il}\,\delta_{jm}-\delta_{im}\,\delta_{jl} .
\]
Por tanto
\[
  [\nab\times(\nab\times\vc{A})]_i=\partial_j\partial_iA_j-\partial_j\partial_jA_i
  =\partial_i\bigl(\partial_jA_j\bigr)-\Lap A_i
  =\bigl[\nab(\nab\cdot\vc{A})\bigr]_i-\bigl[\Lap\vc{A}\bigr]_i ,
\]
donde se han intercambiado $\partial_i$ y $\partial_j$ por Schwarz.
\end{resolucion}

% =========================================================================
\subsection{Seminario 4. Integrales de línea}\label{sem:linea}
% =========================================================================
\begin{ejercicio}[Trabajo a lo largo de una hélice]\label{ejer:helice}
Calcular el trabajo realizado por la fuerza
$\vc{F}(x,y,z)=-\tfrac12x\,\ux-\tfrac12y\,\uy+\tfrac14\,\uz$ al mover una partícula a
lo largo de la hélice $\vc{r}(t)=\cos t\,\ux+\sen t\,\uy+t\,\uz$ desde el punto
$(1,0,0)$ al punto $(-1,0,3\pi)$.
\end{ejercicio}
\begin{resolucion}
\textbf{Paso 1: parametrizar.} En coordenadas, $\gamma(t)=(\cos t,\sen t,t)$ con
$t\in[0,3\pi]$: en $t=0$ estamos en $(1,0,0)$ y en $t=3\pi$, en
$(\cos3\pi,\sen3\pi,3\pi)=(-1,0,3\pi)$. Es decir,
$x(t)=\cos t$, $y(t)=\sen t$, $z(t)=t$.

\textbf{Paso 2: derivar y restringir el campo.} El vector tangente es
$\gamma'(t)=(-\sen t,\cos t,1)$, y sobre la hélice
\[
  \vc{F}(\gamma(t))=\Bigl(-\tfrac12\cos t,\,-\tfrac12\sen t,\,\tfrac14\Bigr).
\]
\textbf{Paso 3: producto escalar.}
\[
  \vc{F}(\gamma(t))\cdot\gamma'(t)=\tfrac12\sen t\cos t-\tfrac12\sen t\cos t+\tfrac14=\tfrac14 .
\]
\textbf{Paso 4: integrar.}
\[
  W=\int_{C}\vc{F}\cdot\dl=\int_{0}^{3\pi}\frac14\,\dif t=\boxed{\frac{3\pi}{4}} .
\]
\end{resolucion}

\begin{observacion}[Comprobación cualitativa]
Los términos con senos y cosenos se han cancelado, y es lo esperable: la parte
horizontal del campo, $-\tfrac12(x\,\ux+y\,\uy)$, es radial (apunta hacia el eje),
mientras que la velocidad horizontal de la partícula es tangente a la
circunferencia, perpendicular a la radial; por tanto esa parte del campo no
realiza trabajo. Solo cuenta la componente vertical, $F_z=\tfrac14$ constante,
que actúa a lo largo de un ascenso total $\Delta z=3\pi$:
$W=\tfrac14\cdot3\pi$ (\cref{fig:sem-helice}).
\end{observacion}

\begin{ejercicio}[Integral de línea sobre un segmento]\label{ejer:segmento}
Calcular $\displaystyle\int_C(x^{2}-y+3z)\,\dif x$, donde $C$ es el segmento
rectilíneo desde el origen hasta $(1,2,1)$.
\end{ejercicio}
\begin{resolucion}
Usamos la parametrización estándar del segmento (\cref{ej:segmento}), con
$\vc{p}=(0,0,0)$ y $\vc{q}=(1,2,1)$:
$\sigma(t)=(1-t)(0,0,0)+t\,(1,2,1)=(t,2t,t)$, $t\in[0,1]$, $\sigma'(t)=(1,2,1)$.
En notación diferencial,
\[
  x=t\Rightarrow\dif x=\dif t,\qquad y=2t\Rightarrow\dif y=2\,\dif t,\qquad
  z=t\Rightarrow\dif z=\dif t .
\]
La 1-forma $(x^{2}-y+3z)\,\dif x$ solo tiene componente en $\dif x$ (equivale a
integrar el campo $(x^{2}-y+3z,\,0,\,0)$), de modo que
\[
  \int_C(x^{2}-y+3z)\,\dif x=\int_{0}^{1}\bigl(t^{2}-2t+3t\bigr)\dif t
  =\int_{0}^{1}(t^{2}+t)\,\dif t=\Bigl[\frac{t^{3}}{3}+\frac{t^{2}}{2}\Bigr]_{0}^{1}
  =\frac13+\frac12=\boxed{\frac56}.
\]
\end{resolucion}

\begin{ejercicio}[Una 1-forma sobre una circunferencia]\label{ejer:circulo}
Calcular $\displaystyle\int_Cx\,\dif y-y\,\dif x$, donde $C$ es la circunferencia
$x^{2}+y^{2}=R^{2}$ recorrida en sentido antihorario.
\end{ejercicio}
\begin{resolucion}
Parametrizamos con el ángulo polar, que es el parámetro natural:
$x=R\cos\varphi$, $y=R\sen\varphi$, $\varphi\in[0,2\pi]$. Entonces
$\dif x=-R\sen\varphi\,\dif\varphi$ y $\dif y=R\cos\varphi\,\dif\varphi$, y componiendo
directamente sobre el integrando diferencial,
\[
  x\,\dif y-y\,\dif x=(R\cos\varphi)(R\cos\varphi)\,\dif\varphi-(R\sen\varphi)(-R\sen\varphi)\,\dif\varphi
  =R^{2}(\cos^{2}\varphi+\sen^{2}\varphi)\,\dif\varphi=R^{2}\,\dif\varphi ,
\]
de donde
\[
  \int_Cx\,\dif y-y\,\dif x=\int_{0}^{2\pi}R^{2}\,\dif\varphi=R^{2}\bigl[\varphi\bigr]_{0}^{2\pi}=\boxed{2\pi R^{2}} .
\]
\end{resolucion}

\begin{observacion}[Interpretación geométrica]
La 1-forma $x\,\dif y-y\,\dif x$ es la integral de línea del campo de giro
$\vc{F}=-y\,\ux+x\,\uy$ (\cref{fig:campos_rot}), que sobre la circunferencia es
tangente, de módulo $R$ y en el sentido de recorrido: $\vc{F}\cdot\dl=R\cdot R\,\dif\varphi$.
Por eso la circulación es $R\cdot2\pi R=2\pi R^{2}$: el doble del área del círculo
($=\pi R^{2}$), en acuerdo con el teorema de Green, pues
$\nab\times\vc{F}=2\,\uz$ (\cref{ej:rot}) y $\oint\vc{F}\cdot\dl=\iint2\,\dif S=2\pi R^{2}$
(\cref{fig:sem-circulo}). Que no se anule indica que $\vc{F}$ \emph{no} es un
campo gradiente (\cref{cor:gradiente}).
\end{observacion}

% Figuras: trabajo sobre una hélice y circulación sobre una circunferencia
\begin{figure}[ht]
\centering
\tdplotsetmaincoords{72}{125}
\begin{tikzpicture}[tdplot_main_coords, scale=1.35]
  \draw[eje] (0,0,0) -- (1.9,0,0) node[below left] {\small $x$};
  \draw[eje] (0,0,0) -- (0,1.9,0) node[right] {\small $y$};
  \draw[eje] (0,0,0) -- (0,0,3.3) node[above] {\small $z$};
  % hélice (eje z comprimido: se dibuja 0.28 t en lugar de t)
  \draw[trayecto, domain=0:{3*pi}, samples=140, variable=\t, smooth]
    plot ({cos(deg(\t))},{sin(deg(\t))},{0.28*\t});
  \foreach \k in {0,1,...,6}{%
    \pgfmathsetmacro{\tt}{\k*pi/2*1.0}%
    \pgfmathsetmacro{\px}{cos(deg(\tt))}\pgfmathsetmacro{\py}{sin(deg(\tt))}\pgfmathsetmacro{\pz}{0.28*\tt}%
    % tangente (azul) y campo (ámbar/carmín)
    \draw[->, line width=1pt, azulprof] (\px,\py,\pz) -- ++({-0.55*sin(deg(\tt))},{0.55*cos(deg(\tt))},0.16);
    \draw[->, line width=1pt, cm3] (\px,\py,\pz) -- ++({-0.42*\px},{-0.42*\py},0.30);
    \fill (\px,\py,\pz) circle (1pt);}
  \fill[rojoteo] (1,0,0) circle (2pt) node[below right, text=black] {\small$(1,0,0)$};
  \fill[rojoteo] (-1,0,{0.28*3*pi}) circle (2pt) node[above left, text=black, fill=white, fill opacity=0.8, text opacity=1, inner sep=1pt] {\small$(-1,0,3\pi)$};
  \node[font=\small, text=azulprof, anchor=west] at (1.55,0.35,2.7) {$\gamma'(t)$};
  \node[font=\small, text=cm3, anchor=west] at (1.55,-0.75,2.2) {$\vc{F}$};
  \node[font=\small, text=rojoteo, anchor=west] at (-0.2,1.5,0.7) {$C$};
\end{tikzpicture}
\caption{Hélice del \cref{ejer:helice}. La parte horizontal de
$\vc{F}=-\tfrac12(x\,\ux+y\,\uy)+\tfrac14\,\uz$ es radial y, por tanto,
perpendicular a la velocidad horizontal $\gamma'$ (tangente a la circunferencia):
solo la componente vertical constante $F_z=\tfrac14$ trabaja. (El eje $z$ se
dibuja comprimido.)}
\label{fig:sem-helice}
\end{figure}

\begin{figure}[ht]
\centering
\begin{tikzpicture}[x=1.0cm, y=1.0cm]
  \shade[inner color=azulclaro!25, outer color=azulclaro!60] (0,0) circle (2.2);
  \clipD{-2.9}{2.9}{-2.9}{2.9}{\flechas{-\y}{\x}{-2.0,-1.5,...,2.0}{-2.0,-1.5,...,2.0}{0.36}{2.9}}
  \draw[trayecto] (0,0) circle (2.2);
  \foreach \a in {20,110,200,290}{%
    \draw[-{Stealth[length=7pt]}, rojoteo, line width=1.6pt] ({2.2*cos(\a)},{2.2*sin(\a)}) -- ++({-0.01*sin(\a)*40},{0.01*cos(\a)*40});}
  \draw[->, line width=0.5pt, grispizarra] (0,0) -- (-1.55,-1.55) node[midway, below left, font=\small] {$R$};
  \node[font=\small, text=rojoteo, anchor=west] at (2.3,1.9) {$C$};
  \node[font=\small, text=azulprof] at (0,0.55) {$\nab\times\vc{F}=2\,\uz$};
  \marcoD{-2.9}{2.9}{-2.9}{2.9}{-2,-1,0,1,2}{-2,-1,0,1,2}
  \ejesD{-2.9}{2.9}{-2.9}{2.9}
\end{tikzpicture}
\caption{Circunferencia del \cref{ejer:circulo}. Sobre $C$, el campo de giro
$\vc{F}=-y\,\ux+x\,\uy$ es tangente, de módulo $R$ y en el sentido de recorrido,
de modo que $\oint\vc{F}\cdot\dl=R\cdot2\pi R=2\pi R^{2}$, el doble del área del
disco, coherente con $\nab\times\vc{F}=2\,\uz$.}
\label{fig:sem-circulo}
\end{figure}


% =========================================================================
\subsection{Seminario 5. Verificación del teorema del gradiente}\label{sem:gradiente}
% =========================================================================
\begin{ejercicio}[Dos caminos, un mismo resultado]\label{ejer:verif-gradiente}
Comprobar que se cumple el teorema del gradiente,
$\int_{\vc{a}}^{\vc{b}}\nab f\cdot\dl=f(\vc{b})-f(\vc{a})$, para el campo escalar
$f:\R^{3}\to\R$ definido por $f(x,y,z)=xy^{2}$ y los puntos
$\vc{a}=(0,0,0)$ y $\vc{b}=(2,2,2)$. Hacer la integral por dos caminos
distintos y comprobar que el resultado es el mismo para ambos.
\end{ejercicio}
\begin{resolucion}
\textbf{El gradiente y el lado derecho del teorema.} El gradiente de $f=xy^{2}$ es
\[
  \nab f=\pd{f}{x}\,\ux+\pd{f}{y}\,\uy+\pd{f}{z}\,\uz=y^{2}\,\ux+2xy\,\uy+0\,\uz .
\]
La diferencia de potencial entre los extremos es
$f(\vc{b})-f(\vc{a})=f(2,2,2)-f(0,0,0)=(2)(2^{2})-0=8$. Este es el número que
tienen que reproducir las dos integrales.

\textbf{Camino 1: segmento rectilíneo $\gamma_1$.} La parametrización estándar
del segmento $[\vc{a},\vc{b}]$ es
$\gamma_1(t)=(1-t)\vc{a}+t\vc{b}=(2t,2t,2t)$, $t\in[0,1]$, de modo que
$\dl_1=\gamma_1'(t)\,\dif t=(2,2,2)\,\dif t$. Evaluamos el gradiente sobre la curva:
\[
  \nab f(\gamma_1(t))=\bigl((2t)^{2},\,2(2t)(2t),\,0\bigr)=(4t^{2},\,8t^{2},\,0).
\]
El producto escalar es $\nab f(\gamma_1)\cdot\gamma_1'=4t^{2}\cdot2+8t^{2}\cdot2+0\cdot2=24t^{2}$, y
\[
  \int_{\gamma_1}\nab f\cdot\dl=\int_{0}^{1}24t^{2}\,\dif t=\Bigl[\frac{24}{3}t^{3}\Bigr]_{0}^{1}=8 .
\]
Coincide con $f(\vc{b})-f(\vc{a})$, como asegura el teorema.

\textbf{Camino 2: ruta ortogonal por los ejes $\gamma_2$.} Ahora elegimos una
trayectoria escalonada, paralela a los ejes cartesianos y dividida en tres
subsegmentos $C_x$, $C_y$, $C_z$ (\cref{fig:sem-trayectorias}), y sumamos sus
contribuciones por aditividad (\cref{prop:prop-linea}):
$\gamma_2=C_x\uplus C_y\uplus C_z$.
\begin{itemize}
  \item \emph{Tramo $C_x$ (recorrido en $x$)}, de $\vc{a}=(0,0,0)$ a $(2,0,0)$:
        $\gamma_{2,x}(t)=(2t,0,0)$, $t\in[0,1]$, $\dl_x=(2,0,0)\,\dif t$. Aquí
        $y=0$, de modo que $\nab f=(y^{2},2xy,0)=(0,0,0)$ y
        $\int_{C_x}\nab f\cdot\dl_x=0$.
  \item \emph{Tramo $C_y$ (recorrido en $y$)}, de $(2,0,0)$ a $(2,2,0)$:
        $\gamma_{2,y}(t)=(2,t,0)$, $t\in[0,2]$, $\dl_y=(0,1,0)\,\dif t$. El
        gradiente sobre el tramo es $\nab f=(t^{2},\,2\cdot2\cdot t,\,0)=(t^{2},4t,0)$,
        y el producto escalar con $\dl_y$ vale $4t\,\dif t$:
        \[
          \int_{C_y}\nab f\cdot\dl_y=\int_{0}^{2}4t\,\dif t=\bigl[2t^{2}\bigr]_{0}^{2}=8 .
        \]
  \item \emph{Tramo $C_z$ (recorrido en $z$)}, de $(2,2,0)$ a $(2,2,2)$:
        $\gamma_{2,z}(t)=(2,2,t)$, $t\in[0,2]$, $\dl_z=(0,0,1)\,\dif t$. El gradiente
        es $\nab f=(4,8,0)$, perpendicular a $\dl_z$:
        $\nab f\cdot\dl_z=0$ y $\int_{C_z}\nab f\cdot\dl_z=0$.
\end{itemize}
Sumando, $\int_{\gamma_2}\nab f\cdot\dl=0+8+0=8=f(\vc{b})-f(\vc{a})$.
\end{resolucion}

\begin{observacion}[Qué hemos comprobado]
Las dos integrales dan $8$ aunque los caminos son muy distintos: es la
\emph{independencia del camino} (\cref{cor:gradiente}). Los tramos tienen
sentido físico: $f=xy^{2}$ no depende de $z$, así que el tramo vertical no
cambia $f$; y mientras $y=0$ la función vale $0$ y no cambia, por lo que todo el
incremento $f:0\to8$ ocurre en el tramo $C_y$, donde $f=2y^{2}$ pasa de $0$ a $8$. Además,
al recorrer $\gamma_1$ y volver por $\gamma_2$ en sentido inverso se obtiene un
circuito cerrado con circulación $8-8=0$ (\cref{eq:circulacion-nula}).
\end{observacion}

% Figura: verificación del teorema del gradiente por dos caminos
\begin{figure}[ht]
\centering
\tdplotsetmaincoords{62}{112}
\begin{tikzpicture}[tdplot_main_coords, scale=1.9]
  \draw[eje] (0,0,0) -- (2.8,0,0) node[below left] {\small $x$};
  \draw[eje] (0,0,0) -- (0,2.8,0) node[right] {\small $y$};
  \draw[eje] (0,0,0) -- (0,0,2.8) node[above] {\small $z$};
  % cota de referencia: aristas del cubo
  \draw[punteado] (2,0,0) -- (2,2,0) -- (0,2,0) (2,2,0) -- (2,2,2) (2,0,0) -- (2,0,2) -- (2,2,2) (0,2,0) -- (0,2,2) -- (2,2,2) (0,0,2) -- (2,0,2) (0,0,2) -- (0,2,2);
  % camino 1: segmento recto
  \draw[line width=1.6pt, rojoteo] (0,0,0) -- (2,2,2);
  \draw[->, line width=1.6pt, rojoteo] (0.9,0.9,0.9) -- (1.05,1.05,1.05);
  % camino 2: poligonal por los ejes
  \draw[line width=1.6pt, azulprof] (0,0,0) -- (2,0,0) -- (2,2,0) -- (2,2,2);
  \draw[->, line width=1.6pt, azulprof] (0.9,0,0) -- (1.1,0,0);
  \draw[->, line width=1.6pt, azulprof] (2,0.9,0) -- (2,1.1,0);
  \draw[->, line width=1.6pt, azulprof] (2,2,0.9) -- (2,2,1.1);
  % gradiente grad f = (y^2, 2xy, 0) sobre C_y y sobre el camino 1
  \foreach \t in {0.4,0.8,1.2,1.6}{%
    \pgfmathsetmacro{\gx}{\t*\t}\pgfmathsetmacro{\gy}{4*\t}\pgfmathsetmacro{\gn}{sqrt(\gx*\gx+\gy*\gy)}
    \mapcolor{\gn/6.6}%
    \draw[-{Stealth[length=4pt]}, line width=0.9pt, mapc] (2,\t,0) -- ++({0.09*\gx},{0.09*\gy},0);}
  \foreach \t in {0.4,0.8,1.2,1.6}{%
    \pgfmathsetmacro{\gx}{\t*\t}\pgfmathsetmacro{\gy}{2*\t*\t}\pgfmathsetmacro{\gn}{sqrt(\gx*\gx+\gy*\gy)}
    \mapcolor{\gn/6.6}%
    \draw[-{Stealth[length=4pt]}, line width=0.9pt, mapc] (\t,\t,\t) -- ++({0.09*\gx},{0.09*\gy},0);}
  \fill[azulprof] (0,0,0) circle (2.2pt) node[below left, text=black] {\small $\vc{a}=(0,0,0)$};
  \fill[azulprof] (2,2,2) circle (2.2pt) node[above right, text=black] {\small $\vc{b}=(2,2,2)$};
  \node[font=\small, text=rojoteo, anchor=east] at (0.85,0.55,1.05) {$\gamma_1$};
  \node[font=\small, text=azulprof] at (1,-0.25,0) {$C_x$};
  \node[font=\small, text=azulprof, anchor=west] at (2.15,1.1,0) {$C_y$};
  \node[font=\small, text=azulprof, anchor=west] at (2.12,2,1.0) {$C_z$};
\end{tikzpicture}
\caption{Verificación del teorema del gradiente para $f=xy^{2}$ entre
$\vc{a}=(0,0,0)$ y $\vc{b}=(2,2,2)$. En rojo, el segmento $\gamma_1$; en azul, la
ruta escalonada $\gamma_2=C_x\uplus C_y\uplus C_z$. Se dibuja también
$\nab f=(y^{2},2xy,0)$ (color según su módulo): en $C_x$ es nulo y en $C_z$ es
perpendicular al camino, de modo que el incremento $f=0\to8$ lo aporta solo
$C_y$. Ambos caminos dan $\int\nab f\cdot\dl=8$.}
\label{fig:sem-trayectorias}
\end{figure}


% =========================================================================
\subsection{Seminario 6. Divergencia y teorema de Gauss}\label{sem:gauss}
% =========================================================================
\begin{ejercicio}[Divergencia de $R^{n}\uR$]\label{ejer:div-Rn}
Demostrar que $\nab\cdot(R^{n}\,\uR)=(n+2)\,R^{n-1}$ para $R\neq0$
(con $n\neq-2$ para que el resultado no sea idénticamente nulo), donde
$\vc{R}=\vc{r}-\vc{r}\,'$.
\end{ejercicio}
\begin{resolucion}
\textbf{Idea.} Ya lo demostramos con reglas de derivación del producto
(\cref{prop:div-rot-Rn}); ahora lo hacemos a mano, componente a componente,
para ver el mecanismo.

Escribimos $R^{n}\uR=R^{n-1}\vc{R}=R^{n-1}\bigl(R_x\,\ux+R_y\,\uy+R_z\,\uz\bigr)$ con
$R_x=x-x'$, etc. La divergencia es la suma de tres derivadas parciales; la primera
requiere la regla del producto:
\[
  \pd{}{x}\bigl[R^{n-1}(x-x')\bigr]=(x-x')\,\pd{R^{n-1}}{x}+R^{n-1}\,\pd{(x-x')}{x} .
\]
Como $\partial R/\partial x=(x-x')/R$ (\cref{eq:dR}), la regla de la cadena da
\[
  \pd{R^{n-1}}{x}=\pd{R^{n-1}}{R}\,\pd{R}{x}=(n-1)R^{n-2}\,\frac{x-x'}{R}
  =(n-1)\,R^{n-3}(x-x') .
\]
Por tanto
$\pd{}{x}\bigl[R^{n-1}(x-x')\bigr]=(n-1)R^{n-3}(x-x')^{2}+R^{n-1}$, y análogamente para
$y$ y $z$:
\[
  \pd{}{y}\bigl[R^{n-1}(y-y')\bigr]=(n-1)R^{n-3}(y-y')^{2}+R^{n-1},\qquad
  \pd{}{z}\bigl[R^{n-1}(z-z')\bigr]=(n-1)R^{n-3}(z-z')^{2}+R^{n-1}.
\]
Sumando los tres, y usando $(x-x')^{2}+(y-y')^{2}+(z-z')^{2}=R^{2}$,
\[
  \nab\cdot\bigl(R^{n}\uR\bigr)=(n-1)\,R^{n-3}R^{2}+3R^{n-1}=(n-1)R^{n-1}+3R^{n-1}=\boxed{(n+2)\,R^{n-1}} .
\]
Este resultado vale para $R\neq0$, es decir, $\vc{r}\neq\vc{r}\,'$. Para $n=-2$,
$\nab\cdot(R^{-2}\uR)=(-2+2)R^{-3}=0$: la divergencia de $\uR/R^{2}$ es nula en todo
punto distinto de la fuente (con la contribución singular de \cref{teo:delta}
en $R=0$).
\end{resolucion}

\begin{ejercicio}[Jacobiano de las coordenadas esféricas]\label{ejer:jacobiano}
Calcular el determinante jacobiano $\partial(x,y,z)/\partial(r,\theta,\varphi)$ del
cambio a coordenadas esféricas.
\end{ejercicio}
\begin{resolucion}
Con $x=r\sen\theta\cos\varphi$, $y=r\sen\theta\sen\varphi$, $z=r\cos\theta$, la
matriz de las derivadas parciales respecto de $(r,\theta,\varphi)$ es
\[
  \frac{\partial(x,y,z)}{\partial(r,\theta,\varphi)}=
  \begin{pmatrix}
    \sen\theta\cos\varphi & r\cos\theta\cos\varphi & -r\sen\theta\sen\varphi\\
    \sen\theta\sen\varphi & r\cos\theta\sen\varphi & \ r\sen\theta\cos\varphi\\
    \cos\theta & -r\sen\theta & 0
  \end{pmatrix}.
\]
Desarrollamos el determinante por la tercera fila (la que tiene un cero):
\begin{align*}
  \det&=\cos\theta\begin{vmatrix}r\cos\theta\cos\varphi & -r\sen\theta\sen\varphi\\ r\cos\theta\sen\varphi & r\sen\theta\cos\varphi\end{vmatrix}
  -(-r\sen\theta)\begin{vmatrix}\sen\theta\cos\varphi & -r\sen\theta\sen\varphi\\ \sen\theta\sen\varphi & r\sen\theta\cos\varphi\end{vmatrix}+0.
\end{align*}
El primer menor vale
$r^{2}\cos\theta\sen\theta\cos^{2}\varphi+r^{2}\cos\theta\sen\theta\sen^{2}\varphi=r^{2}\sen\theta\cos\theta$,
y el segundo
$r\sen^{2}\theta\cos^{2}\varphi+r\sen^{2}\theta\sen^{2}\varphi=r\sen^{2}\theta$
(en ambos, $\cos^{2}\varphi+\sen^{2}\varphi=1$). Entonces
\[
  \det=\cos\theta\,(r^{2}\sen\theta\cos\theta)+r\sen\theta\,(r\sen^{2}\theta)
  =r^{2}\sen\theta\,(\cos^{2}\theta+\sen^{2}\theta)=r^{2}\sen\theta .
\]
El jacobiano (valor absoluto) es $J=|r^{2}\sen\theta|=r^{2}\sen\theta$, porque
$\theta\in[0,\pi]$ implica $\sen\theta\ge0$. Concuerda con el producto de
factores de escala $h_1h_2h_3=1\cdot r\cdot r\sen\theta$ (\cref{prop:elementos}).
\end{resolucion}

\begin{ejercicio}[Teorema de Gauss en una esfera]\label{ejer:gauss-esfera}
Comprobar que se cumple el teorema de la divergencia,
$\int_V\nab\cdot\vc{A}\,\dV=\oint_S\vc{A}\cdot\dS$, siendo $\vc{A}(\vc{r})=r^{2}\,\ur$
(en coordenadas esféricas) y $V$ una esfera de radio $R$ centrada en el origen.
\end{ejercicio}
\begin{resolucion}
\textbf{Idea.} El campo es radial, crece como $r^{2}$ y «sale» de la esfera por
todas partes. Calculamos por separado el lado de volumen y el lado de superficie
y comprobamos que coinciden.

\textbf{Lado del volumen.} La divergencia del campo $\vc{A}=(A_r,A_\theta,A_\varphi)=(r^{2},0,0)$ en esféricas
(\cref{eq:ops-esf}) solo conserva el primer término, pues $A_\theta=A_\varphi=0$:
\[
  \nab\cdot\vc{A}=\frac{1}{r^{2}}\pd{}{r}\bigl(r^{2}A_r\bigr)=\frac{1}{r^{2}}\pd{}{r}\bigl(r^{4}\bigr)=\frac{4r^{3}}{r^{2}}=4r .
\]
La esfera sólida es
$V=\{(x,y,z):x^{2}+y^{2}+z^{2}\le R^{2}\}=\{(r,\theta,\varphi):0\le r\le R,\ 0\le\theta\le\pi,\ 0\le\varphi\le2\pi\}$,
y con el elemento de volumen $\dV=r^{2}\sen\theta\,\dif r\,\dif\theta\,\dif\varphi$
(\cref{ejer:jacobiano}),
\[
  \iiint_V\nab\cdot\vc{A}\,\dV=\int_{0}^{2\pi}\!\!\int_{0}^{\pi}\!\!\int_{0}^{R}(4r)\,(r^{2}\sen\theta)\,\dif r\,\dif\theta\,\dif\varphi .
\]
El integrando es un producto de funciones de una sola variable y los límites son
constantes: por el teorema de Fubini, la integral se factoriza,
\begin{align*}
  &=\Bigl(\int_0^{2\pi}\dif\varphi\Bigr)\Bigl(\int_0^{\pi}\sen\theta\,\dif\theta\Bigr)\Bigl(\int_0^{R}4r^{3}\,\dif r\Bigr)
  =\bigl[\varphi\bigr]_0^{2\pi}\cdot\bigl[-\cos\theta\bigr]_0^{\pi}\cdot\bigl[r^{4}\bigr]_0^{R}\\
  &=2\pi\cdot\bigl[-\cos\pi+\cos0\bigr]\cdot R^{4}=2\pi\cdot2\cdot R^{4}=4\pi R^{4}.
\end{align*}

\textbf{Lado de la superficie.} Sobre la esfera cerrada $S=\partial V$ el radio es
constante, $r=R$, y el campo es $\vc{A}=R^{2}\ur$, de módulo constante y radial. La
normal unitaria exterior de la esfera es $\un=\ur$ (el radio normalizado), de modo
que $\vc{A}\cdot\un=R^{2}\,\ur\cdot\ur=R^{2}$ (producto de dos versores de una base
ortonormal: $\langle\uh{i},\uh{j}\rangle=\delta_{ij}$). El elemento de área en esféricas a $r=R$ es
$\dsc=R^{2}\sen\theta\,\dif\theta\,\dif\varphi$ (\cref{sec:esfericas}), por lo que
\[
  \oiint_S\vc{A}\cdot\un\,\dsc=\int_0^{2\pi}\!\!\int_0^{\pi}R^{2}\cdot R^{2}\sen\theta\,\dif\theta\,\dif\varphi
  =R^{4}\Bigl(\int_0^{2\pi}\dif\varphi\Bigr)\Bigl(\int_0^{\pi}\sen\theta\,\dif\theta\Bigr)
  =R^{4}\cdot2\pi\cdot2=4\pi R^{4}.
\]

\textbf{Conclusión.} Ambos lados valen $4\pi R^{4}$, así que
$\iiint_V\nab\cdot\vc{A}\,\dV=\oiint_S\vc{A}\cdot\dS=4\pi R^{4}$ y el teorema de Gauss se
verifica para este campo radial y geometría esférica.
\end{resolucion}

\begin{observacion}
Es instructivo ver por qué sale así. El campo crece hacia fuera ($A_r=r^{2}$) y
cada capa esférica «deja pasar» más flujo hacia fuera de la que recibe de
dentro; esa diferencia es lo que mide la divergencia ($4r>0$ en todo punto,
fuente distribuida). La suma de todas esas contribuciones infinitesimales,
ponderadas por el volumen de cada capa, reproduce el flujo total que atraviesa
la superficie exterior (\cref{fig:sem-esfera}).
\end{observacion}

% Figura: verificación del teorema de la divergencia en una esfera
\begin{figure}[ht]
\centering
\begin{minipage}[t]{0.52\textwidth}\centering
\tdplotsetmaincoords{70}{125}
\begin{tikzpicture}[tdplot_main_coords, scale=1.05, baseline=(current bounding box.north)]
  \def\RR{2.1}
  \draw[eje] (0,0,0) -- (3.0,0,0) node[below left] {\small $x$};
  \draw[eje] (0,0,0) -- (0,3.1,0) node[right] {\small $y$};
  \draw[eje] (0,0,0) -- (0,0,3.1) node[above] {\small $z$};
  % esfera: contorno aparente y círculos máximos
  \begin{scope}[tdplot_screen_coords]
    \shade[ball color=azulclaro!55, opacity=0.35] (0,0) circle (\RR);
    \draw[azulprof!60, line width=0.6pt] (0,0) circle (\RR);
  \end{scope}
  \draw[azulprof!45, line width=0.4pt, domain=0:360, samples=72] plot ({\RR*cos(\x)},{\RR*sin(\x)},0);
  \foreach \ph in {0,45,90,135}{\draw[azulprof!35, line width=0.4pt, domain=0:360, samples=72]
     plot ({\RR*sin(\x)*cos(\ph)},{\RR*sin(\x)*sin(\ph)},{\RR*cos(\x)});}
  % campo radial A = R^2 u_r (flechas salientes, módulo constante)
  \foreach \th/\ph in {35/20,35/200,70/110,70/290,110/20,110/200,145/110,145/290}{%
    \pgfmathsetmacro{\px}{\RR*sin(\th)*cos(\ph)}\pgfmathsetmacro{\py}{\RR*sin(\th)*sin(\ph)}\pgfmathsetmacro{\pz}{\RR*cos(\th)}%
    \draw[-{Stealth[length=4pt]}, line width=1pt, acento] (\px,\py,\pz) -- ++({0.55*sin(\th)*cos(\ph)},{0.55*sin(\th)*sin(\ph)},{0.55*cos(\th)});}
  % elemento de superficie y normal
  \def\ta{48}\def\tb{62}\def\pa{28}\def\pb{46}
  \fill[ocreteo, opacity=0.7]
    ({\RR*sin(\ta)*cos(\pa)},{\RR*sin(\ta)*sin(\pa)},{\RR*cos(\ta)}) --
    ({\RR*sin(\tb)*cos(\pa)},{\RR*sin(\tb)*sin(\pa)},{\RR*cos(\tb)}) --
    ({\RR*sin(\tb)*cos(\pb)},{\RR*sin(\tb)*sin(\pb)},{\RR*cos(\tb)}) --
    ({\RR*sin(\ta)*cos(\pb)},{\RR*sin(\ta)*sin(\pb)},{\RR*cos(\ta)}) -- cycle;
  \draw[-{Stealth[length=5pt]}, line width=1.5pt, azulacero]
    ({\RR*sin(55)*cos(37)},{\RR*sin(55)*sin(37)},{\RR*cos(55)}) --
    ({(\RR+1.0)*sin(55)*cos(37)},{(\RR+1.0)*sin(55)*sin(37)},{(\RR+1.0)*cos(55)});
  \node[font=\small, text=azulacero, anchor=south west] at ({(\RR+1.0)*sin(55)*cos(37)},{(\RR+1.0)*sin(55)*sin(37)},{(\RR+1.0)*cos(55)}) {$\dS$};
  \fill (0,0,0) circle (1.2pt);
\end{tikzpicture}\\[2pt]{\small\scshape (a)}
\end{minipage}\hfill
\begin{minipage}[t]{0.44\textwidth}\centering
\begin{tikzpicture}[baseline=(current bounding box.north)]
  \begin{axis}[tema, width=\textwidth, height=6.3cm, xlabel={$r$}, xmin=0, xmax=1.6, ymin=0, ymax=42,
      samples=100, domain=0:1.6, xtick={0,0.5,1,1.5}, xticklabels={0,0{,}5,$R$,1{,}5},
      xtick={0,0.5,1.0,1.5}, legend style={at={(0.03,0.97)}, anchor=north west}]
    \addplot[azulacero, line width=1.1pt] {4*pi*x^4};
    \addplot[acento, line width=1.1pt] {16*pi*x^3};
    \addplot[name path=curva, acento, draw=none, domain=0:1] {16*pi*x^3};
    \path[name path=eje] (axis cs:0,0) -- (axis cs:1,0);
    \addplot[acento!25] fill between[of=curva and eje];
    \addplot[grispizarra, densely dashed] coordinates {(1,0) (1,12.57)};
    \node[font=\scriptsize, text=azulacero, anchor=west] at (axis cs:1.05,9.5) {$4\pi R^{4}$};
    \legend{$\Phi(r)=4\pi r^{4}$, $16\pi r^{3}=\nab\cdot\vc{A}\,\cdot4\pi r^{2}$}
  \end{axis}
\end{tikzpicture}\\[2pt]{\small\scshape (b)}
\end{minipage}
\caption{Teorema de Gauss para $\vc{A}=r^{2}\,\ur$ en una esfera de radio
$R$. (a)~El campo es radial y de módulo constante sobre la esfera; la normal
exterior es $\ur$, de modo que $\vc{A}\cdot\un=R^{2}$ y el flujo es
$R^{2}\cdot4\pi R^{2}$. (b)~Visión «capa a capa»: cada cáscara de radio $r$ y
espesor $\dif r$ aporta $\nab\cdot\vc{A}\,\dV=4r\cdot4\pi r^{2}\,\dif r$; el área
bajo $16\pi r^{3}$ (sombreada) hasta $R=1$ es el flujo total
$\Phi(R)=4\pi R^{4}$.}
\label{fig:sem-esfera}
\end{figure}


% =========================================================================
\subsection{Seminario 7. Verificación del teorema de Stokes}\label{sem:stokes}
% =========================================================================
\begin{ejercicio}[Stokes en un triángulo]\label{ejer:stokes-triangulo}
Comprobar que se cumple el teorema del rotacional,
$\iint_S(\nab\times\vc{A})\cdot\dS=\oint_{\partial S}\vc{A}\cdot\dl$, siendo
$\vc{A}(\vc{r})=xy\,\ux+2yz\,\uy+3zx\,\uz$ y $S$ el área encerrada por el triángulo
rectángulo de vértices $\vc{a}=(0,0,0)$, $\vc{b}=(2,0,0)$ y $\vc{c}=(0,2,0)$ (en el
plano $z=0$).
\end{ejercicio}
\begin{resolucion}
\textbf{Idea.} Calculamos el flujo del rotacional por el triángulo y la
circulación del campo por sus tres lados, y comprobamos que coinciden. Antes de
empezar fijamos la orientación: recorremos el borde $L=\partial S$ en el orden
$\vc{a}\to\vc{b}\to\vc{c}\to\vc{a}$ (sentido antihorario visto desde $+z$); por la
regla de la mano derecha, la normal compatible es $+\uz$ (\cref{fig:sem-stokes}).

\textbf{1. Integral de superficie.} \emph{El rotacional:}
\[
  \nab\times\vc{A}=
  \begin{vmatrix}
    \ux & \uy & \uz\\
    \partial_x & \partial_y & \partial_z\\
    xy & 2yz & 3zx
  \end{vmatrix}
  =(0-2y)\,\ux-(3z-0)\,\uy+(0-x)\,\uz=-2y\,\ux-3z\,\uy-x\,\uz .
\]
\emph{La superficie:} $S$ está en el plano $z=0$; la parametrizamos con
$\sigma(x,y)=(x,y,0)$, con $(x,y)\in S^{*}$ la región plana limitada por los catetos y la
hipotenusa $y=-x+2$:
\[
  S^{*}=\{(x,y):0\le x\le2,\ 0\le y\le-x+2\}.
\]
Las derivadas parciales son $\partial_x\sigma=(1,0,0)=\ux$ y
$\partial_y\sigma=(0,1,0)=\uy$, y el vector normal asociado es
$\vc{n}_\sigma=\partial_x\sigma\times\partial_y\sigma=\ux\times\uy=\uz$, compatible con la orientación
elegida. \emph{El integrando:} evaluado sobre la superficie ($z=0$),
$(\nab\times\vc{A})(\sigma(x,y))=(-2y,\,0,\,-x)$, y su producto escalar con la normal
es $(\nab\times\vc{A})\cdot\uz=-x$. Por tanto
\[
  \iint_S(\nab\times\vc{A})\cdot\dS=\int_0^2\!\!\int_0^{-x+2}(-x)\,\dif y\,\dif x
  =\int_0^2-x\bigl[y\bigr]_0^{-x+2}\dif x=\int_0^2(x^{2}-2x)\,\dif x
  =\Bigl[\frac{x^{3}}{3}-x^{2}\Bigr]_0^2=\frac83-4=-\frac43 .
\]

\textbf{2. Integral de línea (circulación).} Dividimos el contorno en tres tramos,
$L=\Gamma_1\uplus\Gamma_2\uplus\Gamma_3$:
\begin{itemize}
  \item \emph{Tramo $\Gamma_1$} ($\vc{a}\to\vc{b}$, a lo largo del eje $x$): $\gamma_1(t)=(t,0,0)$, $t\in[0,2]$,
        $\dl=(1,0,0)\,\dif t$. Aquí $y=z=0$, luego $\vc{A}=(xy,2yz,3zx)=\vc{0}$ y $\int_{\Gamma_1}\vc{A}\cdot\dl=0$.
  \item \emph{Tramo $\Gamma_2$} ($\vc{b}\to\vc{c}$, la hipotenusa): $z=0$, $y=-x+2$, y por
        tanto $\dif y=-\dif x$. Parametrizamos con $t\in[0,2]$ mediante
        $\gamma_2(t)=(2-t,\,t,\,0)$ (es el segmento $(1-s)\vc{b}+s\,\vc{c}$ con $s=t/2$), de modo que
        $\gamma_2'(t)=(-1,1,0)$ y $\dl=(-1,1,0)\,\dif t$. Sobre el tramo, $\vc{A}=\bigl((2-t)t,\,0,\,0\bigr)$ (porque $z=0$), y
        $\vc{A}\cdot\gamma_2'=-(2-t)t=t^{2}-2t$, así que
        \[
          \int_{\Gamma_2}\vc{A}\cdot\dl=\int_0^2(t^{2}-2t)\,\dif t=\Bigl[\frac{t^{3}}{3}-t^{2}\Bigr]_0^2=\frac83-4=-\frac43 .
        \]
  \item \emph{Tramo $\Gamma_3$} ($\vc{c}\to\vc{a}$, a lo largo del eje $y$): $\gamma_3(t)=(0,2-t,0)$, $t\in[0,2]$,
        $\dl=(0,-1,0)\,\dif t$. Aquí $x=z=0$, luego $\vc{A}=\vc{0}$ y la integral es nula.
\end{itemize}
Sumando,
\[
  \oint_{\partial S}\vc{A}\cdot\dl=0-\frac43+0=-\frac43 .
\]

\textbf{Conclusión.} Ambos lados valen $-4/3$:
\[
  \iint_S(\nab\times\vc{A})\cdot\dS=\oint_{\partial S}\vc{A}\cdot\dl=-\frac43 ,
\]
y el teorema de Stokes se verifica formalmente.
\end{resolucion}

\begin{observacion}[Lectura geométrica]
Solo el lado $\Gamma_2$ contribuye a la circulación, pues sobre los ejes
coordenados ($y=0$ o $x=0$, con $z=0$) el campo se anula. En el interior del
triángulo, la componente que cuenta del rotacional, $(\nab\times\vc{A})\cdot\uz=-x$, es
negativa y más intensa hacia el vértice $\vc{b}$ (donde $x$ es mayor) (\cref{fig:sem-stokes}):
el campo tiende a hacer girar una rueda en sentido \emph{horario} visto desde
$+z$, justo al contrario del sentido en que recorremos el borde, por eso el
resultado es negativo.
\end{observacion}

% Figura: verificación del teorema de Stokes en un triángulo
\begin{figure}[ht]
\centering
\begin{minipage}[t]{0.48\textwidth}\centering
\begin{tikzpicture}[>=Stealth, x=1.9cm, y=1.9cm, baseline=(current bounding box.north)]
  \draw[eje] (-0.25,0) -- (2.55,0) node[right, text=black] {$x$};
  \draw[eje] (0,-0.25) -- (0,2.55) node[above, text=black] {$y$};
  \fill[azulclaro!55] (0,0) -- (2,0) -- (0,2) -- cycle;
  \draw[trayecto] (0,0) -- (2,0) -- (0,2) -- cycle;
  \draw[-{Stealth[length=8pt]}, rojoteo, line width=1.6pt] (0.85,0) -- (1.15,0);
  \draw[-{Stealth[length=8pt]}, rojoteo, line width=1.6pt] (1.07,0.93) -- (0.93,1.07);
  \draw[-{Stealth[length=8pt]}, rojoteo, line width=1.6pt] (0,1.15) -- (0,0.85);
  \fill[azulprof] (0,0) circle (1.8pt) node[below left, text=black] {$\vc{a}$};
  \fill[azulprof] (2,0) circle (1.8pt) node[below right, text=black] {$\vc{b}$};
  \fill[azulprof] (0,2) circle (1.8pt) node[above left, text=black] {$\vc{c}$};
  \node[font=\small, text=rojoteo, anchor=north] at (1,-0.05) {$\Gamma_1$};
  \node[font=\small, text=rojoteo, anchor=south west] at (1.05,1.05) {$\Gamma_2:\ y=-x+2$};
  \node[font=\small, text=rojoteo, anchor=east] at (-0.06,1) {$\Gamma_3$};
  \node[font=\small, text=azulprof] at (0.48,0.42) {$S$};
  % normal saliendo del plano
  \draw[azulacero, line width=0.8pt] (0.62,0.62) circle (0.075);
  \fill[azulacero] (0.62,0.62) circle (0.02);
  \node[font=\footnotesize, text=azulacero, anchor=south west] at (0.66,0.64) {$\un=\uz$};
  \node[font=\scriptsize, text=black!75, anchor=north west] at (0.15,-0.28) {$\vc{A}=\vc{0}$ sobre $\Gamma_1$ y $\Gamma_3$};
\end{tikzpicture}\\[2pt]{\small\scshape (a)}\ contorno y orientación
\end{minipage}\hfill
\begin{minipage}[t]{0.48\textwidth}\centering
\begin{tikzpicture}[>=Stealth, x=1.9cm, y=1.9cm, baseline=(current bounding box.north)]
  \foreach \i in {0,...,79}{%
    \pgfmathsetmacro{\xa}{2*\i/80}\pgfmathsetmacro{\xb}{2*(\i+1)/80}%
    \pgfmathsetmacro{\tt}{(\xa+\xb)/4}\mapcolor{\tt}%
    \fill[mapc, draw=mapc, line width=0.2pt] (\xa,0) -- (\xb,0) -- (\xb,{2-\xb}) -- (\xa,{2-\xa}) -- cycle;}
  \draw[eje] (-0.25,0) -- (2.55,0) node[right, text=black] {$x$};
  \draw[eje] (0,-0.25) -- (0,2.55) node[above, text=black] {$y$};
  \draw[black!60, line width=0.8pt] (0,0) -- (2,0) -- (0,2) -- cycle;
  % escala de color
  \foreach \i in {0,...,29}{%
    \pgfmathsetmacro{\yy}{0.2+1.4*\i/30}\pgfmathsetmacro{\tt}{\i/30}\mapcolor{\tt}%
    \fill[mapc] (2.35,\yy) rectangle (2.5,{\yy+1.4/30});}
  \draw[black!60, line width=0.4pt] (2.35,0.2) rectangle (2.5,1.6);
  \node[font=\scriptsize, anchor=west] at (2.5,0.2) {$0$};
  \node[font=\scriptsize, anchor=west] at (2.5,0.9) {$-1$};
  \node[font=\scriptsize, anchor=west] at (2.5,1.6) {$-2$};
  \node[font=\scriptsize, anchor=south] at (2.42,1.68) {$(\nab\times\vc{A})\cdot\uz$};
  \node[font=\small, text=white, fill=cm0, fill opacity=0.65, text opacity=1, inner sep=2pt] at (0.62,0.55) {$\iint_S=-\tfrac43$};
\end{tikzpicture}\\[2pt]{\small\scshape (b)}\ integrando $(\nab\times\vc{A})\cdot\uz=-x$
\end{minipage}
\caption{Verificación del teorema de Stokes para $\vc{A}=xy\,\ux+2yz\,\uy+3zx\,\uz$ sobre el
triángulo de vértices $\vc{a},\vc{b},\vc{c}$. (a)~El borde se recorre
$\vc{a}\to\vc{b}\to\vc{c}\to\vc{a}$ (antihorario visto desde $+z$) y, por la
regla de la mano derecha, la normal es $\uz$. El campo se anula sobre los ejes y
solo $\Gamma_2$ contribuye a la circulación. (b)~La componente normal del
rotacional $-x$ es más intensa (en valor absoluto) hacia el vértice $\vc{b}$; su
integral sobre el triángulo es $-4/3$, igual a la circulación $-4/3$.}
\label{fig:sem-stokes}
\end{figure}

